\section{Identified comparisons and a nonlinear economic counterfactual}
\label{sec:r39evidence}

The original exact-construction cases remain below, including the faster
structural comparator and the earlier adverse reuse comparisons. They are not
used to identify a precision effect. We instead analyze the deposited R38
factorial and nonlinear execution, and add an independently checked all-state
network certificate. No historical clock is replaced by a reconstruction.

\subsection{Design and accounting}
The fixed protocol contains 315 isolated services: 288 quadratic policy
services, 12 spline services, 12 trained-network services and three nonlinear
accuracy-frontier services. There are 105 deterministic economy--method objects
with three timing repetitions each. A separate warm-up is excluded from the
contrasts. Execution order is shuffled with a fixed seed, and the numerical
libraries use one thread. Neither CPU affinity nor frequency was controlled.
The displayed time is the complete child service through output synchronization;
the archive also records startup-inclusive parent time. Minima and maxima
accompany medians rather than a significance claim from three observations.
Memory is process high-water RSS, not exclusive algorithm memory.

The quadratic design crosses fixed binary64 versus adaptive native precision
with inverse rebuilding versus within-target caching. It adds a cached fixed
method selected from binary32 and binary64, in that order, and the structural
solution. All tuning trials, failed policy checks and precision retries enter
the service boundary. The selected fixed type is the lowest successful type in
that two-element set, not the minimum possible arbitrary mantissa length.
The update rule is held fixed in this factorial; comparisons with the older
update rules remain separately identified in the historical table.

Every candidate policy is checked initially and after an actor change; the
service stops at its first successful economic certificate. The target grid is
$10^{-2},10^{-4},10^{-6},10^{-8},10^{-10}$. Further objects vary $(d,T)$ over
$(8,4),(16,8),(32,4),(4,12)$ with condition parameters $8,16,16,4$, change the
valuation by factors $1.005,1/4,4$, include cold and transported starts, and
include an additional moment-margin case. This finite design tests several
operating regimes; it is not an asymptotic scaling experiment or an exhaustive
stress test near the entropic moment boundary.

\begin{table}[htbp]
\centering\small\setlength{\tabcolsep}{5pt}
\caption{Precision and caching at a common policy target}\label{tab:r39factorial}
\begin{tabular}{lrrrrr}
\toprule
Method & Median (s) & Range (s) & Updates & Builds & Loss bound\\
\midrule
Fixed64, rebuilt & 0.416 & [0.416, 0.420] & 11 & 11 & $5.59\times10^{-6}$\\
Fixed64, cached & 0.415 & [0.415, 0.421] & 11 & 3 & $5.59\times10^{-6}$\\
Adaptive, rebuilt & 0.421 & [0.418, 0.422] & 11 & 11 & $5.59\times10^{-6}$\\
Adaptive, cached & 0.433 & [0.417, 0.433] & 11 & 3 & $5.59\times10^{-6}$\\
Tuned native, cached & 0.424 & [0.423, 0.428] & 11 & 3 & $5.59\times10^{-6}$\\
Structural & 0.170 & [0.168, 0.173] & 0 & 0 & $4.46\times10^{-28}$\\
\bottomrule
\end{tabular}
\par\smallskip\begin{minipage}{.96\linewidth}\footnotesize $d=2$, $T=4$, target $10^{-4}$. Three isolated complete-service clocks per object. Builds count target inverses, not actor solves. All recorded failed checks and tuning are charged.\end{minipage}
\end{table}


The factorial removes the original inverse-caching confound. It does not show
a stable timing advantage from adaptive precision. At tolerance $10^{-4}$,
the fixed64 cached median is about $0.415$ seconds and the adaptive cached
median about $0.433$ seconds; the structural median is about $0.170$ seconds.
The two cached factor procedures each make eleven updates and three inverse
constructions. The uncached procedures construct eleven inverses. Reusing an
inverse changes this count, but verification and own-policy evaluation can
dominate the small-matrix service clock.

\begin{table}[htbp]
\centering\small\setlength{\tabcolsep}{5pt}
\caption{First certified policy across accuracy targets}\label{tab:r39accuracy}
\begin{tabular}{rrrrrr}
\toprule
Target & Adaptive bound & Bound/target & Adaptive (s) & Fixed64 (s) & Structural (s)\\
\midrule
0.01 & 0.0023 & 0.23 & 0.352 & 0.353 & 0.175\\
0.0001 & $5.59\times10^{-6}$ & 0.0559 & 0.433 & 0.415 & 0.170\\
$1\times10^{-6}$ & $4.68\times10^{-11}$ & $4.68\times10^{-5}$ & 0.444 & 0.442 & 0.171\\
$1\times10^{-8}$ & $4.35\times10^{-11}$ & 0.00435 & 0.500 & 0.482 & 0.169\\
$1\times10^{-10}$ & $4.35\times10^{-11}$ & 0.435 & 0.486 & 0.491 & 0.175\\
\bottomrule
\end{tabular}
\par\smallskip\begin{minipage}{.96\linewidth}\footnotesize Factor methods use caching. Times are medians; complete dispersion and all methods appear in the supplement. Discrete oversolving is displayed, not suppressed.\end{minipage}
\end{table}


First-pass stopping removes unnecessary continuation after a successful policy
check. It does not force a discrete update to land at the target. The achieved
bound-to-target ratio for adaptive caching ranges from approximately
$4.68\times10^{-5}$ to $0.435$ over these five targets. In particular, the
$10^{-6}$ object still oversolves substantially. Table~\ref{tab:r39accuracy}
therefore reports achieved error as well as work; it is not presented as a
frontier with uniformly binding constraints. The structural solver also
oversolves, for a different reason: the economy permits an exact structural
policy. All 288 policy services certify their declared targets. Complete
counters, process memory and timing dispersion for every cell are supplied in
the supplement and the machine-readable audit.

Small and large transported changes have different operational consequences.
For adaptive caching, the $1.005$ change uses three new factor updates rather
than twelve from a cold start. The factor-$4$ change fails all three transported
factor gates and falls back to cold construction. The factor-$1/4$ change also
fails the transported gates. Setup of the old policy is charged jointly in the
reported warm service; accordingly its total clock need not beat a cold solve.
These are observed gate outcomes, not a success rate estimated for an unknown
population of future changes.

\subsection{Nonlinear approximation, policy accuracy and economic gains}
The economy in Section~\ref{sec:r39economy} is run for four decision dates at
$(P,\theta)\in\{1,4\}\times\{0,1\}$. The direct construction uses $N=256$
state cells and $A=128$ action cells. The trained continuation has 31 hidden
units, changes all affine hidden weights, and uses 300 Adam steps per date.
Training data are deterministic certified Bellman nodal values, not independent
observations. Every state and every continuous feasible action are covered by
the interpolation and action allowances in the theorem. The finite shock sum
is evaluated through interval operations with Taylor enclosures for exponential
and logarithm; no analytic Gaussian continuation is used.

\begin{table}[htbp]
\centering\small\setlength{\tabcolsep}{5pt}
\caption{Nonlinear capital: policy bounds and complete work}\label{tab:r39nonlinear}
\begin{tabular}{rrrrrr}
\toprule
Price & Risk & Neural bound & Spline bound & Neural (s) & Spline (s)\\
\midrule
1 & 0 & 0.2359 & 0.2352 & 2.846 & 0.240\\
1 & 1 & 0.2366 & 0.2359 & 5.011 & 1.311\\
4 & 0 & 0.2533 & 0.2520 & 2.999 & 0.409\\
4 & 1 & 0.2542 & 0.2528 & 6.222 & 2.564\\
\bottomrule
\end{tabular}
\par\smallskip\begin{minipage}{.96\linewidth}\footnotesize $T=4$, $N=256$, $A=128$. Cost includes reference construction and, for neural services, all fitting and verification. Price-four services also include construction and evaluation of the old rule. Bounds are in model cost units.\end{minipage}
\end{table}


The trained policies have global loss bounds between approximately $0.236$ and
$0.254$. These are deliberately reported in the model's cost units, not
reinterpreted as $10^{-4}$ certificates. The spline method is faster in all
four complete-service comparisons. A separate constructive accuracy sweep
has first-pass bounds $0.9352,0.4704,0.2359,0.1181$ at targets
$1,1/2,1/4,1/8$, respectively, using $N=64,128,256,512$ and $A=N/2$.
The corresponding bound-to-target ratios lie between $0.935$ and $0.946$.
Its cumulative construction clocks include the unsuccessful coarser meshes;
they are reported separately from the full service clock and are not four
independent experimental observations from each repetition.

The price change nevertheless supports a substantive policy comparison.
The old rule is the certified conventional rule constructed at $P=1$ and then
held fixed; the new rule is selected by the trained neural continuation at
$P=4$. Both are evaluated under the new price, by an independent finite-shock
tree enclosure. Table~\ref{tab:r39economic} reports the second bracket in
\eqref{eq:r39decomposition}. Its lower bound is positive at each reported
initial state for both risk specifications. Thus the new policy strictly
improves on retaining the old rule under the changed investment burden.
This comparison is not a claim that the neural rule outperforms the newly
optimized spline rule. The latter remains the stronger work comparator.

\begin{table}[htbp]
\centering\small\setlength{\tabcolsep}{5pt}
\caption{Cost saved by reoptimization after the investment-price change}\label{tab:r39economic}
\begin{tabular}{rrrrrr}
\toprule
Risk & Capital & Old action & New action & Gain lower & Gain upper\\
\midrule
0 & 0.125 & 0.25000 & 0.21094 & 0.085570 & 0.085571\\
0 & 0.25 & 0.25000 & 0.16992 & 0.097610 & 0.099736\\
0 & 0.5 & 0.16797 & 0.10938 & 0.057079 & 0.057080\\
0 & 0.75 & 0.06836 & 0.05273 & 0.021237 & 0.030832\\
1 & 0.125 & 0.25000 & 0.21289 & 0.084553 & 0.084554\\
1 & 0.25 & 0.25000 & 0.17383 & 0.097370 & 0.098241\\
1 & 0.5 & 0.16797 & 0.10938 & 0.057067 & 0.057068\\
1 & 0.75 & 0.07031 & 0.05273 & 0.023443 & 0.028670\\
\bottomrule
\end{tabular}
\par\smallskip\begin{minipage}{.96\linewidth}\footnotesize Old: conventional rule constructed at price one, held fixed at price four. New: neural-selected rule at price four. Gains compare both rules under the new price. Displayed gain endpoints are rounded outward to six decimals; action entries are descriptive.\end{minipage}
\end{table}


Higher investment burden is accompanied here by lower initial investment in
low and intermediate capital states. The fixed-rule exposure formula explains
why keeping the old action rule incurs a positive repricing burden, while the
independent policy comparison quantifies how much of that burden is avoided.
The gain is economically interpretable even though the uniform global
optimality bound is more conservative. It does not identify an empirical
elasticity or a calibrated welfare effect.

\subsection{New all-state network verification}
The endpoint verifier in Proposition~\ref{prop:r39breakpoints} is applied to
all 60 deposited date-specific network records. Deduplication by coefficients
and reference function gives 17 distinct networks, not 60 independent fits.
Every new uniform network-to-spline bound is below $0.001$; the previous
mesh-Lipschitz checks were between $0.027$ and $0.031$. Independent exact
rational evaluation of every affine-piece endpoint verifies all 17 bounds,
with two further fixtures covering zero slopes and nondyadic roots.
This is an improvement in verification resolution, not a new timing experiment
or an automatic improvement in the previously recorded policy-loss numbers.
The audit records exact numerators and denominators, outward upper bounds,
source identities and the mapping from every service to its distinct network.

\paragraph{Interpretation.}
The results establish a constructive NBO comparison outside exact quadratics,
a residual-verified native factor implementation, and certified nonlinear
reoptimization gains. They also show why a certificate is not a work-superiority
result. The structural quadratic solver and the nonlinear spline baseline
remain favorable comparators. General diffusion, endogenous-preference,
temporal-self and game applications retain their original statements and
application-specific obligations; the two completed constructions are not
silently substituted for those distinct assumptions.
